% ============================================================
%  CAPÍTULO 2: MARCO INSTITUCIONAL Y CONTEXTO SANITARIO
% ============================================================
\chapter{Marco institucional y contexto sanitario en España}
\label{cap:contexto}

% ------------------------------------------------------------
\section{Introducción}
% ------------------------------------------------------------

Comprender el comportamiento de los hogares españoles respecto a la sanidad
privada exige situar su decisión en el contexto institucional específico del
Sistema Nacional de Salud. España no dispone de un único sistema sanitario, sino
de \textbf{diecisiete sistemas autonómicos} con diferencias sustanciales en
dotación de recursos, modelos de gestión y resultados en salud. Esta heterogeneidad
institucional es precisamente la que permite el análisis econométrico de panel que
desarrollamos en este trabajo. 

El presente capítulo describe la arquitectura del SNS, analiza la evolución de las
listas de espera desde 2016, examina el papel del seguro privado como complemento
del sistema público y justifica la pertinencia del enfoque autonómico para el
estudio del fenómeno.

% ------------------------------------------------------------
\section{Organización del Sistema Nacional de Salud}
% ------------------------------------------------------------

\subsection{Los principios fundacionales}

El Sistema Nacional de Salud español fue establecido por la Ley General de
Sanidad de 1986, inspirada en el modelo Beveridge británico. Sus principios
rectores son la \textbf{universalidad} (cobertura de toda la población), la
\textbf{equidad} (acceso independiente de la capacidad de pago), la
\textbf{financiación pública} (mediante impuestos generales) y la
\textbf{integralidad} de las prestaciones (prevención, diagnóstico, tratamiento y
rehabilitación).

Estos principios configuran un sistema en el que, teóricamente, ningún ciudadano
debería necesitar acudir al sector privado por motivos económicos o de acceso. La
realidad, sin embargo, muestra un crecimiento sostenido de los seguros privados
que sugiere que el sistema público no satisface completamente las expectativas de
una parte significativa de la población.

\subsection{La descentralización sanitaria}

El proceso de transferencia de competencias sanitarias a las comunidades autónomas,
culminado en 2002, transformó radicalmente la gobernanza del SNS. Actualmente, cada
CCAA dispone de un \textbf{Servicio de Salud propio} con competencias plenas en
planificación, gestión y provisión de servicios sanitarios. El Ministerio de
Sanidad conserva funciones de coordinación, política farmacéutica, sanidad exterior
y garantía de la cohesión del sistema, pero carece de autoridad ejecutiva sobre los
servicios regionales.

Esta arquitectura ha dado lugar a lo que algunos autores denominan \textit{``17
sistemas de salud''} que, compartiendo principios comunes, difieren en:

\begin{itemize}
  \item \textbf{Gasto sanitario per cápita:} En 2022, el gasto público por
    habitante oscilaba entre 1.200~€ (Andalucía) y 1.900~€ (País Vasco), un rango
    del 60\%.
  \item \textbf{Ratio de profesionales:} Las diferencias en médicos y enfermeros
    por 1.000 habitantes son sustanciales, con comunidades que superan la media
    europea y otras que se sitúan por debajo.
  \item \textbf{Modelos de gestión:} Algunas CCAA han experimentado con fórmulas
    de colaboración público-privada, externalizaciones o conciertos, mientras que
    otras mantienen modelos de gestión directa integral.
  \item \textbf{Infraestructura hospitalaria:} La densidad de camas hospitalarias,
    equipamiento tecnológico y accesibilidad geográfica varía significativamente.
\end{itemize}

\subsection{Implicaciones para las listas de espera}

La descentralización ha generado una \textbf{competencia implícita en eficiencia}
entre comunidades autónomas. Los tiempos de espera, al ser un indicador visible y
políticamente sensible, se han convertido en un campo de comparación permanente.
Los informes semestrales del Ministerio de Sanidad sobre listas de espera reciben
amplia cobertura mediática, y los gobiernos autonómicos se ven presionados a
mejorar sus resultados relativos.

Sin embargo, esta competencia opera con restricciones presupuestarias
heterogéneas. Las comunidades con mayor capacidad fiscal (Madrid, País Vasco,
Navarra, Cataluña) disponen de márgenes de maniobra superiores para aumentar la
oferta y reducir tiempos de espera. Las comunidades con menor renta per cápita
enfrentan una tensión estructural entre demanda sanitaria creciente (a menudo
asociada a mayor envejecimiento) y recursos limitados.

% ------------------------------------------------------------
\section{Evolución de las listas de espera (2016--2024)}
% ------------------------------------------------------------

\subsection{Tendencia general al alza}

El análisis de los datos del Sistema de Información sobre Listas de Espera del
SNS (SISLE-SNS) revela una tendencia estructural al \textbf{deterioro progresivo}
de los tiempos de espera desde 2016. La espera media para consulta de especialistas
pasó de 58 días en diciembre de 2016 a 95 días en diciembre de 2023, un incremento
del 64\% en siete años.

Esta tendencia no ha sido uniforme: se observan períodos de relativa estabilización
(2017--2018), episodios de deterioro acelerado (2019, 2022--2023) y el fenómeno
anómalo de la pandemia (2020--2021). Pero la dirección general es inequívoca:
\textbf{el sistema público es progresivamente menos capaz de atender la demanda en
tiempos razonables}.

\subsection{Factores estructurales: el envejecimiento demográfico}

El principal factor explicativo de esta tendencia es el \textbf{envejecimiento de
la población}. España tiene una de las estructuras demográficas más envejecidas
del mundo, con un 20{,}5\% de población mayor de 65 años (media del panel) y
comunidades como Castilla y León o Asturias que superan el 25\%.

El envejecimiento impacta doblemente en las listas de espera:

\begin{itemize}
  \item \textbf{Por el lado de la demanda:} Los mayores consumen servicios
    sanitarios con mucha mayor intensidad que los jóvenes. La prevalencia de
    enfermedades crónicas, pluripatología y necesidad de intervenciones quirúrgicas
    (cataratas, prótesis articulares, cirugía cardíaca) aumenta exponencialmente
    con la edad.
  \item \textbf{Por el lado de la oferta:} Una población más envejecida implica
    una base fiscal más reducida (menor tasa de empleo, menores ingresos medios),
    lo que limita la capacidad de financiación del sistema público.
\end{itemize}

\subsection{El \textit{shock} de la pandemia COVID-19}

La pandemia de 2020 supuso un \textbf{\textit{shock} exógeno} de magnitud sin
precedentes para el sistema sanitario español. Durante los meses más críticos
(marzo--mayo de 2020), la actividad programada se suspendió casi por completo para
concentrar recursos en la atención a pacientes COVID. Las listas de espera
quirúrgicas se vaciaron artificialmente ---no porque se operase a los pacientes,
sino porque se dejó de añadir nuevos---, para después acumular una demanda
embalsada que explotó en 2021--2022.

El efecto neto de la pandemia sobre los tiempos de espera fue paradójico:

\begin{itemize}
  \item \textbf{Corto plazo (2020):} Aparente reducción de listas (suspensión de
    derivaciones).
  \item \textbf{Medio plazo (2021--2023):} Explosión de tiempos de espera al
    reactivarse la demanda acumulada.
  \item \textbf{Largo plazo:} Revelación de las fragilidades estructurales del
    sistema y aceleración de las reformas en algunas CCAA.
\end{itemize}

\subsection{Heterogeneidad autonómica}

La Tabla~\ref{tab:espera_ccaa} presenta los tiempos de espera en consulta de
especialistas para las CCAA extremas del panel en diciembre de 2023.

\begin{table}[htbp]
  \centering
  \caption{Tiempos de espera para consulta de especialista por CCAA (diciembre 2023)}
  \label{tab:espera_ccaa}
  \begin{threeparttable}
    \begin{tabular}{lrc}
      \toprule
      \textbf{CCAA} & \textbf{Espera media (días)} & \textbf{Posición} \\
      \midrule
      País Vasco         &  42 & Mejor \\
      Madrid             &  51 & 2.ª mejor \\
      Navarra            &  54 & 3.ª mejor \\
      \multicolumn{3}{c}{\vdots} \\
      Cataluña           &  98 & 14.ª \\
      Canarias           & 124 & 15.ª \\
      Castilla-La Mancha & 147 & 16.ª \\
      Extremadura        & 156 & Peor \\
      \bottomrule
    \end{tabular}
    \begin{tablenotes}[flushleft]
      \footnotesize
      \item Fuente: Ministerio de Sanidad, SISLE-SNS, diciembre 2023.
    \end{tablenotes}
  \end{threeparttable}
\end{table}

El rango observado ---de 42 a 156 días--- implica que un ciudadano de Extremadura
espera casi \textbf{cuatro veces más} que uno del País Vasco para acceder a la
misma prestación nominal del SNS. Esta heterogeneidad es la que justifica el
análisis a nivel autonómico y convierte las CCAA en \textit{``laboratorios
naturales''} para estudiar el efecto de las listas de espera sobre la demanda de
seguros privados.

% ------------------------------------------------------------
\section{El sector privado: definición y papel en el sistema sanitario}
% ------------------------------------------------------------

\subsection{Modalidades de seguro sanitario privado}

En España, el seguro sanitario privado puede clasificarse en dos grandes categorías:

\textbf{Seguros individuales:} Contratados directamente por el ciudadano con una
compañía aseguradora, cubren al titular y opcionalmente a su unidad familiar. El
mercado está dominado por compañías como Sanitas, Adeslas, Asisa, DKV y MAPFRE
Salud.

\textbf{Seguros colectivos:} Contratados por empresas como beneficio para sus
empleados, ofrecen condiciones ventajosas (primas reducidas por economías de
escala, eliminación de cuestionarios médicos). Este segmento ha crecido
significativamente en la última década al incorporarse la cobertura sanitaria como
elemento de retribución flexible y retención de talento.

Ambas modalidades comparten la característica de ser \textbf{complementarias} al
SNS, no sustitutivas: el asegurado privado mantiene su derecho a la sanidad
pública y puede utilizar ambos sistemas según conveniencia.

\subsection{Dimensión del sector privado en España}

Según datos de ICEA, a finales de 2023:

\begin{itemize}
  \item Aproximadamente \textbf{12 millones de españoles} (25\% de la población)
    disponían de algún tipo de seguro sanitario privado.
  \item El volumen de primas del ramo de salud superó los \textbf{11.500 millones
    de euros}, con un crecimiento interanual del 6{,}2\%.
  \item La ratio de siniestralidad se situó en torno al 82\%.
\end{itemize}

La penetración del seguro privado no es homogénea territorialmente: supera el
35\% en Madrid y Cataluña, mientras que en comunidades como Extremadura,
Castilla-La Mancha o Galicia se sitúa por debajo del 15\%.

\subsection{El seguro privado como ``válvula de escape''}

Desde la perspectiva del análisis económico, el seguro sanitario privado funciona
como una \textbf{válvula de escape} del sistema público. Cuando las listas de
espera del SNS se alargan más allá de lo tolerable para ciertos hogares, estos
pueden optar por contratar un seguro que les garantice acceso rápido a consultas,
pruebas diagnósticas e intervenciones. Este mecanismo tiene tres implicaciones
fundamentales:

\begin{description}
  \item[\textbf{Equidad:}] La válvula de escape solo está disponible para quienes
    pueden pagar la prima del seguro. El resultado es un sistema sanitario
    \textbf{formalmente universal pero materialmente estratificado} según capacidad
    de pago.
  \item[\textbf{Eficiencia:}] La existencia de la válvula de escape reduce la
    presión política sobre el sistema público ---el fenómeno \textit{``exit versus
    voice''} de \textcite{Hirschman1970}--- y puede contribuir paradójicamente al
    deterioro del sector público.
  \item[\textbf{Sostenibilidad:}] A medio plazo, la migración de población joven y
    sana hacia el sector privado puede generar \textbf{selección adversa} en el
    SNS.
\end{description}

% ------------------------------------------------------------
\section{Justificación del enfoque autonómico}
% ------------------------------------------------------------

El análisis a nivel \textbf{autonómico} ---y no nacional--- se justifica por
varias razones. Primera, la descentralización sanitaria implica que las condiciones
de acceso al SNS varían sustancialmente entre CCAA; datos agregados nacionales
producirían estimaciones sesgadas. Segunda, la decisión de contratar seguros
privados se toma a nivel del hogar, que enfrenta los tiempos de espera de
\textit{su} servicio de salud autonómico, no una media nacional abstracta. Tercera,
la variación entre CCAA proporciona la \textbf{variación estadística} necesaria
para identificar efectos causales. Cuarta, el análisis de panel permite controlar
por \textbf{efectos fijos autonómicos} y explotar la \textbf{variación temporal
dentro de cada comunidad}.

% ------------------------------------------------------------
\section{Síntesis del contexto institucional}
% ------------------------------------------------------------

El marco institucional descrito configura el escenario en el que los hogares
españoles toman decisiones sobre contratación de seguros sanitarios privados:

\begin{enumerate}
  \item Un sistema público universal pero descentralizado en 17 servicios de salud
    con eficiencias muy dispares.
  \item Una tendencia estructural al deterioro de los tiempos de espera, impulsada
    por el envejecimiento demográfico y agravada por la pandemia.
  \item Un sector privado consolidado que ofrece una alternativa accesible para
    quienes pueden pagar las primas.
  \item Una heterogeneidad territorial que permite identificar el efecto de las
    listas de espera sobre la demanda privada mediante técnicas de panel.
\end{enumerate}

Los capítulos siguientes formalizan este contexto en un marco teórico
(Capítulo~\ref{cap:teoria}), lo operacionalizan con datos
(Capítulo~\ref{cap:datos}), y contrastan las hipótesis mediante estimación
econométrica (Capítulos~\ref{cap:metodologia} y \ref{cap:resultados}).
